\BourbakiStarSection{Bibliography}

\begin{enumerate}
\def\labelenumi{\arabic{enumi}.}
\item
  O. NEUGEBAUER, Vorlesungen über Geschichte der antiken Mathematik, Vol. I: Vorgriechische Mathematik, Berlin (Springer), 1934.
\item
  Euclidis Elementa, 5 vols., ed.~J. L. Heiberg, Lipsiae (Teubner), 1883--88.
\end{enumerate}

2 bis. T. L. HEATH, The thirteen books of Euclid's Elements . . ., 3 vols., Cambridge, 1908.

\begin{enumerate}
\def\labelenumi{\arabic{enumi}.}
\setcounter{enumi}{2}
\item
  FRANCISCI VIETAE, Opera mathematica . . ., Lugduni Batavorum (Elzevir), 1646.
\item
  P. FERMAT, Oeuvres, vol.~I, Paris (Gauthier-Villars), 1891: (a) Ad locos planos et solidos Isagoge (pp.~91--110; French transl. ibid., vol.~III, p.~85); (b) Appendix ad methodum \ldots{} (pp.~184--188; French transl., ibid., vol.~III, p.~159).
\item
  R. DESCARTES, La géométrie, Leyde (Jan Maire), 1637 (= Oeuvres, ed.~Ch. Adam and P. Tannery, vol.~VI, Paris (L. Cerf), 1902).
\item
  G. DESARGUES, Oeuvres . . ., vol.~I, Paris (Leiber), 1864: Brouillon proiect d'une atteinte aux éuénements des rencontres d'un cône avec un plan (pp.~103--230).
\item
  G. W. LEIBNIZ, Mathematische Schriften, ed.~C. I. Gerhardt, vol.~I, Berlin (Asher), 1849.
\item
  L. EULER: (a) Introductio in Analysin Infinitorum, vol.~2, Lausannae, 1748 (= Opera Omnia (1), vol.~IX, Zürich-Leipzig-Berlin (O. Füssli and B. G. Teubner), 1945); (b) Institutionum Calculi Integralis, vol.~2, Petropoli, 1769 (= Opera Omnia (1), vol.~XII, Leipzig-Berlin (B. G. Teubner), 1914).
\item
  J.-L. LAGRANGE, Oeuvres, Paris (Gauthier-Villars), 1867--1892: (a) Solutions analytiques de quelques problèmes sur les pyramides triangulaires,
\end{enumerate}

vol.~III, pp.~661--692; (b) Solution de différents problèmes de calcul intégral, vol.~I, p.~471; (c) Recherches d'arithmétique, vol.~III, pp.~695--795.

\begin{enumerate}
\def\labelenumi{\arabic{enumi}.}
\setcounter{enumi}{9}
\item
  G. CRAMER, Introduction à l'analyse des lignes courbes, Geneva (Cramer and Philibert), 1750.
\item
  E. BEZOUT, Théorie générale des équations algébriques, Paris, 1779.
\item
  C. F. GAUSS, Werke, Göttingen, 1870--1927: (a) Disquisitiones arithmeticae, vol.~I; (b) Selbstanzeige zur Theoria residuorum biquadraticorum, Commentatio secunda, vol.~II, pp.~169--178.
\item
  A.-L. CAUCHY, Mémoire sur les fonctions qui ne peuvent obtenir que deux valeurs égales et de signes contraires par suite des transpositions opérées entre les variables qu'elles renferment, J. Ec. Polytech., cahier 17 (volume X), 1815, pp.~29--112 (= Oeuvres complètes (2), vol.~I, Paris (Gauthier-Villars), 1905, pp.~91--169).
\item
  A.-L. CAUCHY, in Leçons de calcul différentiel et de calcul intégral, rédigées principalement d'après les méthodes de M. A.-L. Cauchy by Abbé Moigno, vol.~II, Paris, 1844.
\item
  P. G. LEJEUNE-DIRICHLET, Werke, vol.~I, Berlin (G. Reimer), 1889, pp.~619--644.
\item
  C. G. J. JACOBI, Gesammelte Werke, Berlin (G. Reimer), 1881--1891: (a) De formatione et proprietatibus determinantium, vol.~III, pp.~355--392; (b) De functionibus duarum variabilium\ldots, vol.~II, pp.~25--50.
\item
  M. CHASLES, Aperçu historique sur l'origine et le développement des méthodes en géométrie \ldots, Bruxelles, 1837.
\item
  A. F. Möbius, Der baryzentrische Calcul\ldots, Leipzig, 1827 (=Gesammelte Werke, vol.~I, Leipzig (Hirzel), 1885).
\item
  H. GRASSMANN: (a) Die lineale Ausdehnungslehre, ein neuer Zweig der Mathematik, dargestellt und durch Anwendungen auf die übrigen Zweige der Mathematik, wie auch auf die Statik, Mechanik, die Lehre vom Magnetismus und die Kristallonomie erläutert, Leipzig (Wigand), 1844 (=Gesammelte Werke, vol.~I, 1st Part, Leipzig (Teubner), 1894); (b) Die Ausdehnungslehre, vollständig und in strenger Form bearbeitet, Berlin, 1862 (=Gesammelte Werke, vol.~I, 2nd Part, Leipzig (Teubner), 1896).
\item
  W. R. HAMILTON, Lectures on quaternions, Dublin, 1853.
\item
  J. J. SYLVESTER, Collected Mathematical Papers, vol.~I, Cambridge, 1904: No.~25, Addition to the articles . . ., pp.~145--151 (=Phil. Mag., 1850).
\item
  A. CAYLEY, Collected Mathematical Papers, Cambridge, 1889--1898: (a) Sur quelques théorèmes de la géométrie de position, vol.~I, pp.~317--328 (=Crelle's J., vol.~XXXI (1846), pp.~213--227); (b) A memoir on the theory of matrices, vol.~II, pp.~475--496 (=Phil. Trans., 1858).
\item
  K. WEIERSTRASS, Mathematische Werke, vol.~II, Berlin, (Mayer und Müller),
\end{enumerate}

1895: Zur Theorie des aus n Haupteinheiten gebildeten complexen Grössen, pp.~311--332.

\begin{enumerate}
\def\labelenumi{\arabic{enumi}.}
\setcounter{enumi}{23}
\item
  R. DEDEKIND, Gesammelte mathematische Werke, 3 vols., Braunschweig (Vieweg), 1930--32.
\item
  H. J. SMITH, Collected Mathematical Papers, vol.~I, Oxford, 1894: On systems of linear indeterminate equations and congruences, p.~367 (=Phil. Trans., 1861).
\item
  L. KRONECKER, Vorlesungen über die Theorie der Determinanten . . ., Leipzig (Teubner), 1903.
\item
  G. PEANO, Calcolo geometrico secondo l'Ausdehnungstheorie di Grassmann preceduto dalle operazioni della logica deduttiva, Torino, 1888.
\item
  G. RICCI and T. LEVI-CIVITA, Méthodes de calcul différentiel absolu et leurs applications, Math. Ann., vol.~LIV (1901), p.~125.
\item
  E. CARTAN, Sur certaines expressions différentielles et le problème de Pfaff, Ann. E.N.S. (3), vol.~XVI (1899), pp.~239--332 (= Oeuvres complètes, vol.~II \(_{1}\) , Paris (Gauthier-Villars), 1953, pp.~303--396).
\item
  H. POINCARÉ, Les méthodes nouvelles de la mécanique céleste, vol.~III, Paris (Gauthier-Villars), 1899, Chapter XXII.
\item
  O. TOEPLITZ, Ueber die Auflösung unendlichvieler linearer Gleichungen mit unendlichvielen Unbekannten, Rend. Circ. Mat. Pal., vol.~XXVIII (1909), pp.~88--96.
\end{enumerate}

